\begin{table}[tp]
\centering
\caption{Necessary oracle budgets for stationarity with \href{https://github.com/k3nfalt/longtails/blob/main/LongTails/Setups/NonnegativeSmoothSampleLossEuclidean.lean}{nonnegative smooth sample losses in Euclidean space}.}\label{tab:NonnegativeSmoothSampleLossEuclidean}
\parbox{\linewidth}{\footnotesize Let $L,\Delta,\varepsilon>0$ and $R=L\Delta/\varepsilon^2\ge2^{64}$. Each $\ell_i\ge0$ has an $L$-Lipschitz gradient, and $F=\mathbb E_i\ell_i$ satisfies $F(0)\le\Delta$. Any randomized adaptive algorithm may query two points per call, receiving only their gradients for the same fresh IID sample. The target is $\mathbb E\|\nabla F(\hat x)\|\le\varepsilon$.}\par\smallskip
\begin{threeparttable}
\footnotesize
\setlength{\tabcolsep}{4pt}\renewcommand{\arraystretch}{1.25}
\begin{tabularx}{\linewidth}{@{}>{\raggedright\arraybackslash}p{0.20\linewidth}>{\raggedright\arraybackslash}p{0.28\linewidth}>{\raggedright\arraybackslash}p{0.28\linewidth}>{\raggedright\arraybackslash}X@{}}
\toprule
\textbf{Result} & \textbf{Oracle complexity} & \textbf{Novelty} & \textbf{Limitations} \\
\midrule
\textbf{(1)} \href{https://github.com/k3nfalt/longtails/blob/main/LongTails/Theorems/NonnegativeSampleStationarityLowerBound.lean}{Nonnegative\allowbreak Sample\allowbreak Stationarity\allowbreak Lower\allowbreak Bound} \textbf{(Ours)} $(\star\star)$ \href{https://github.com/k3nfalt/longtails/blob/main/scripts/reviews/NonnegativeSampleStationarityLowerBound.md}{(review)}
& Lower bound: $s^3+1=\Omega(R^{3/2})$ pair calls. The hard instance uses $m=[512(s+1)]^2$ equiprobable samples and dimension $\Theta(s^{11})$ at minimum width. & & \\
\midrule[1.2pt]
\textbf{(2)} \citep[Thm.~3.1]{pengTangJiaDenseV3}
& Lower bound: $\Omega(m+\sqrt m\,R)$ component calls for $m\ge m_0$, $\varepsilon^2\le c_0L\Delta$ and fixed failure probability $0<\delta<p_0$. Two calls simulate one uniform pair.
& \textbf{(1)} proves $\Omega(R^{3/2})$ for nonnegative losses with bounded initial risk. \textbf{(2)} gives that order at $m\asymp R$ under a gap bound. Whether shifting its components to nonnegativity preserves $F(0)\le\mathcal{O}(\Delta)$ is open.
& \textbf{(1)} ties $m$ to accuracy and hides sample indices and values. \textbf{(2)} permits adaptive component selection and returns component values and gradients. \\
\midrule
\textbf{(3)} \citep[Thm.~3, Eq.~(24)]{hdabcArjevaniV2}
& Lower bound: $\Omega\!\left(\frac{\Delta\bar L\sigma}{\varepsilon^3}+\frac{\sigma^2}{\varepsilon^2}\right)$ same-sample pair calls for $\varepsilon\le c\sqrt{\bar L\Delta}$.
& \textbf{(1)} establishes hardness under individual smoothness, nonnegativity and bounded initial risk. \textbf{(3)} gives the same order formally at $\bar L=L$, $\sigma^2=L\Delta$ for a mean-squared smooth oracle.
& \textbf{(1)} covers gradient-only replies. \textbf{(3)} also allows exact objective values. \\
\bottomrule
\end{tabularx}
\begin{tablenotes}[flushleft]\footnotesize
\item[] \textbf{Abbr:} $s=\lfloor\sqrt R/2^{32}\rfloor$; $m$ is the component count. In \textbf{(2)}, \textbf{(3)}, $\Delta$ bounds $F(0)-\inf F$. In \textbf{(3)}, $\bar L$ is the mean-squared smoothness parameter and $\sigma^2$ bounds the oracle variance. $m_0,c_0,p_0,c$ are universal constants.
\end{tablenotes}
\end{threeparttable}
\end{table}
